% \section{Introduction}\label{sec:Intro}

% Humanitarian crises, such as natural disasters, conflicts, and pandemics, often result in widespread devastation, displacing populations and disrupting essential services. Timely and effective delivery of humanitarian aid is crucial to mitigate the impact of such crises and save lives~\parencite{smith2005}. However, humanitarian logistics operations face numerous challenges, including uncertainty in demand, limited resources, logistical constraints, and dynamic operating environments. Optimal resource allocation plays a central role in addressing these challenges by ensuring that available resources are allocated efficiently to meet the needs of affected populations.
% \begin{equation}
% \mathrm{Cost\_Red}_j=\rho_0+\rho_1N_{m_j}+\rho_2N_{l_j}+\rho_3I2_{m_j}+\rho_4I2_{l_j}+\rho_5OS_{m_j}+\phi\mathbf{X}_j+\mu_j
% \end{equation}

% In recent years, there has been growing interest in applying mathematical optimization techniques to improve the efficiency and effectiveness of humanitarian logistics operations. Traditional approaches typically rely on deterministic models that assume perfect information and static conditions. However, in practice, humanitarian crises are characterized by uncertainty and variability, making it essential to incorporate stochastic elements into the modeling process. Stochastic programming offers a powerful framework for addressing uncertainty and optimizing decision-making under uncertain conditions.
% \begin{align}
% \mbox{Cost\_Red}_j&=\rho_0+\rho_1N_{m_j}+\rho_2N_{l_j}+\rho_3I2_{m_j}\notag\\
% &\quad+\rho_4I2_{l_j}+\rho_5OS_{m_j}+\phi\mathbf{X}_j+\mu_j
% \end{align}

% In this paper, we present a stochastic programming approach to optimize resource allocation in humanitarian logistics~\parencitep{jones2010}. Our approach aims to balance the trade-off between cost minimization and service quality while considering uncertainty in demand, supply, and transportation constraints. We develop a mixed-integer linear programming (MILP) formulation that captures the stochastic nature of humanitarian crises and provides decision support for relief agencies. The proposed model integrates both deterministic and stochastic components, allowing decision-makers to make informed decisions in uncertain environments.

% Text of your paper here
\section{Introduction}
% 1. What is money-laundering fraud?
%Illicit finance represents a profound challenge to the stability of global markets.
Financial institutions operate in an increasingly complex environment where identifying high-risk customer behaviors has become both a regulatory imperative and a strategic business challenge. As AI technologies reshape financial services and redefine business decision-making processes~(\cite{agrawal2019artificial}), financial institutions must balance predictive accuracy with operational constraints, regulatory compliance, and customer relationship value. Between 2020 and 2024, U.S. banks filed more than 137,000 suspicious activity reports linked to illicit networks~\parencite{website:FinCEN}, while recent investigation revealed that criminal organizations moved over \$312 billion through legitimate financial channels ~\parencite{website:WSJ}. Globally, the United Nations estimates that 2–5\% of world GDP (\$800 billion–\$2 trillion) is laundered annually, yet less than 1\% is ever recovered~\parencite{website:UN}. Beyond regulatory penalties, ineffective risk management imposes substantial operational costs through false alerts requiring manual investigation, potential reputation damage from compliance failures, and lost legitimate customer relationships due to over-cautious interventions. Understanding how well-intentioned financial regulations can produce unintended consequences is crucial for designing effective AML policies~\parencite{bao2017ni}.

% 2. Background of money laundering
Traditional customer risk management systems rely predominantly on rule-based heuristics that flag transactions exceeding predetermined thresholds or matching known suspicious patterns. While straightforward to implement, these systems generate overwhelming false positive rates—often exceeding 95\%—creating a critical resource allocation problem: compliance teams must investigate alerts that consume investigator time without commensurate risk reduction, while legitimate customers experience unnecessary friction that erodes relationship value. More critically, rigid rule-based approaches fail to adapt to evolving customer behavioral patterns and cannot detect sophisticated schemes that deliberately fragment activities across multiple institutions, jurisdictions, and time periods. As financial transactions become increasingly digital, interconnected, and global, the limitations of static rule-based customer monitoring have become economically costly and strategically problematic. 

% 3. Emphasize the importance of machine learning methods. Show the advantages of machine learning, including efficiency and effectiveness of dealing with large-scale data. Use the existing money-laundering methods as evidence.
Recent advances in machine learning offer promising alternatives for customer risk assessment. Data-driven approaches can process massive, complex behavioral datasets with efficiency and adaptability that traditional rules cannot match. By automatically identifying latent behavioral patterns and relationship anomalies, these methods can improve both detection accuracy and operational efficiency. However, despite their analytical promise, applying machine learning to customer risk management in financial markets faces four fundamental business and technical challenges that have limited widespread adoption. 

% 4. Introduce the challenges of existing money laundering methods, including the privacy concern of sharing data across banks, interpretability of model, and the lack of decision-making after detection.
The first challenge is the fragmented customer intelligence across competing institutions. Financial institutions operate as competing entities, constrained by privacy regulations, competitive concerns, and customer confidentiality requirements. Customer transaction data are distributed across different entities and platforms, and institutions cannot easily share customer information due to strict regulatory constraints and competitive dynamics~\parencite{amoako2025exploring}. This fragmented view significantly limits the ability to capture illicit behavioral patterns that span multiple institutions—a deliberate strategy employed by sophisticated actors who exploit the lack of cross-institutional coordination. Models trained on isolated institutional data are prone to high false negatives, where suspicious activities appear benign within a single institution's customer base but reveal clear patterns when viewed across the broader market ecosystem. For instance, a transaction series may appear routine within one bank's customer portfolio but reveal layering or structuring patterns only when combined with behavioral data from other institutions. Addressing this challenge requires developing privacy-preserving collaborative frameworks that enable institutions to benefit from shared market intelligence without exposing proprietary customer information or violating competitive boundaries. 

The second challenge is lack of interpretability in customer risk models. Most advanced predictive models, particularly deep learning architectures, operate as "black boxes," making it difficult for relationship managers, compliance officers, and regulators to understand why particular customers or transactions are flagged. In high-stakes financial environments, model transparency is not merely desirable but mandatory for ensuring accountability, supporting customer relationship decisions, and satisfying regulatory auditing standards. When a model flags a corporate customer for suspicious wire transfers, investigators must trace the reasoning—such as anomalous transaction frequency, counterpart network characteristics, or deviations from historical behavioral patterns—to justify the alert and determine appropriate relationship management actions. Without interpretability, analysts face opaque predictions they cannot validate or contest, leading to mistrust and resistance toward model adoption. Moreover, regulators require financial institutions to explain the logic behind automated decisions, especially when these affect customer rights, relationship continuity, or trigger legal actions. The lack of interpretability thus impedes both internal investigation workflows and creates significant compliance and governance risks, ultimately limiting the business value of sophisticated analytics investments. 

Thirdly, detection is merely the first step in an effective customer risk management pipeline, as institutions must decide whether to monitor, escalate, or intervene in customer relationships under operational constraints and time pressures. Current systems typically employ rigid rule-based escalation procedures that fail to align with the probabilistic outputs of predictive models or optimize across multiple business objectives. A flagged transaction involving cross-border remittances might warrant different interventions depending on the customer's relationship value, transaction context, historical behavior patterns, and broader network risk exposure. Without adaptive decision-support mechanisms, relationship managers must manually reconcile probabilistic risk signals with static compliance rules, leading to inconsistent interventions, relationship friction, and suboptimal resource allocation. False positives—incorrectly flagging legitimate customers—can lead to unnecessary account restrictions, customer dissatisfaction, and relationship attrition, while false negatives—failing to identify genuine risks—result in financial losses and severe regulatory penalties. These challenges are compounded by institutional fragmentation, where each institution observes only a local subset of the broader customer relationship network. There is an urgent need for adaptive targeting frameworks that coordinate detection with intervention policies across institutional boundaries while optimizing for relationship value, operational costs, compliance requirements, and customer experience.

The fourth challenge is the extreme behavioral class imbalance. In real-world financial systems, illicit transactions constitute a minute fraction of total transaction volume—often less than 1\% as documented in Table~\ref{tab:country_graph_stats}—while the overwhelming majority corresponds to legitimate customer activities. This extreme imbalance fundamentally challenges model development and deployment. Standard machine learning algorithms become biased toward the majority (legitimate) class, achieving deceptively high overall accuracy while exhibiting poor sensitivity to rare but critical suspicious behaviors. Models may fail to detect subtle anomalies hidden within massive streams of normal customer transactions. Moreover, suspicious behaviors are highly heterogeneous and evolve over time, meaning the limited labeled positive cases cannot adequately capture the diversity of real-world schemes such as layering, structuring, or trade-based techniques. The scarcity of verified suspicious cases also reflects the costly and time-intensive nature of financial investigations, where labeling requires expert validation and cross-institutional cooperation. This imbalance complicates not only model training and evaluation but also operational deployment, as systems learn to overlook minority patterns that deviate from dominant behavioral trends. 



% 5. The introduction of the major component of our proposed method
To address these challenges, this paper proposes a privacy-preserving adaptive framework that connects three critical dimensions of modern financial security: illicit activities detection module, interpretable group identification, and hierarchical decision-making. Collectively, these innovations enable financial institutions to move from reactive detection toward proactive, coordinated, and explainable anti-money-laundering operations. Specifically, we first introduce a graph-based federated learning framework designed to detect money-laundering patterns while preserving data privacy. In this approach, each financial institution analyzes its own transaction network locally but shares abstracted, non-sensitive insights via \textit{virtual supernode} to corporately learn broader patterns of suspicious laundering behavior. This enables institutions to benefit from shared intelligence without exposing proprietary or customer data. The detection module also integrates multiple analytical objectives, such as identifying hidden transaction clusters and predicting anomalous relationships, to better capture the complexity of real-world laundering schemes. Instead of using the cross-entropy loss, we apply focal loss to mitigate the data imbalanced issue by dynamically adjusting the relative importance of the minority (suspicious) versus majority (normal) classes. The results show that the Cross-Entropy Loss tends to favor the majority class (normal transactions), leading to lower false alarm rates (Type I errors) but a higher tendency to overlook true laundering cases (Type II errors). In contrast, the Focal Loss introduces a more balanced learning process by assigning greater weight to these hard-to-detect, high-risk transactions, thus achieving much lower Type II errors and slightly higher Type I errors. We then evaluate the individual contribution of each component (e.g., virtual super-node, graph cut loss) through an ablation study, where components are systematically removed and replaced. The findings demonstrate that every component contributes meaningfully to the overall detection performance.

This paper reveals one fundamental challenge for cross-border financial crime detection: \textit{subpopulation shift} induced by data fragmentation resulting in performance degradation. \textit{subpopulation shift} is a form of distribution shift where the overall population remains fixed, but the characteristics or proportions of underlying subgroups differ~\parencite{koh2021wilds}. When transaction data are partitioned by country to reflect regulatory and privacy constraints, substantial distributional discrepancies emerge across markets, particularly for critical features such as transaction amount, which spans several orders of magnitude. Our experiments demonstrate that commonly used feature normalization strategies can inadvertently exacerbate this issue. Specifically, normalizing features independently within each country amplifies subpopulation shift, as country-level statistics differ markedly from global ones. Controlled experiments comparing country-level and global-level normalization show that this induced shift significantly degrades the performance of most machine learning models, leading to higher false negatives and lower precision-recall performance. Most approaches benefit from eliminating subpopulation shift via globally consistent normalization. These findings highlight subpopulation shift as a critical and previously underappreciated obstacle in AML modeling and motivate the need for learning framework that are robust to fragmented, privacy-constrained and heterogeneous financial data environments.

This paper also highlight several insights central to the effectiveness of detection module. First, federated learning generates substantial value from coordination without data sharing, a property particularly valuable when regulatory constraints prohibit direct information exchange. By aggregating model parameters rather than raw transactions, institutions achieve detection gains equivalent to pooled training while maintaining competitive independence. This reduces both Type I and Type II errors, demonstrating that shared model parameters (not shared data) deliver meaningful gains in detection performance. Second, models that incorporate structural information in graph data, such as federated graph neural network, consistently outperform feature-only deep learning methods, such as three-layer MLP. By leveraging relational signals that describe how accounts interact, the graph-based model better distinguishes legitimate from suspicious behaviors. 
Third, collaborative learning efficiency is critical: varying communication frequency in federated training demonstrates a non-monotonic trade-off between detection accuracy and synchronization cost. Moderate communication intervals improve AUPRC and reduce missed detections while lowering coordination overhead, whereas overly frequent or infrequent aggregation can degrade performance. 
% Fourth, our analysis reveals that data fragmentation itself harms detection quality: when traditional machine learning models are trained separately in each country, their performance systematically declines compared with training on aggregated data. This finding underscores the importance of collaborative learning across institutions and motivates the federated framework, which achieves the benefits of unified training while maintaining strict privacy requirements. 
Fourth, data fragmentation and subpopulation shift negatively impact detection quality: normalizing features independently within each country or training models solely on country-level subsets reduces performance compared with using aggregated or globally normalized data. To address this, we design a fixed-value min–max normalization strategy to mitigate subpopulation shift without sharing sensitive raw data across institutions. Overall, these findings underscore the importance of collaborative, privacy-preserving learning frameworks that can reconcile performance, structural modeling, and data heterogeneity in cross-country anti–money laundering detection.


Upon these detection results by the detection module, we further employ a network propagation mechanism via cross-bank Personalized PageRank (PPR) that starts from one suspicious account and expands the identification of suspicious neighbors entities, mimicking the process of money flows. This technique reveals sub-group structures within suspicious transaction networks, essentially identifying clusters of accounts or entities that act in coordination. By uncovering these latent community patterns, the module provides analysts and regulators with clearer explanations of what type of money-laundering behaviors are flagged, helping bridge the gap between algorithmic output and human reasoning. 
% By leveraging the group suspicious money-laundering behaviors, the subsequent label propagation process further refines these insights, spreading information across the transaction network to detect indirect or secondary participants that might otherwise remain hidden. 
Together, these mechanisms turn abstract detection scores into interpretable group-level evidence, empowering financial analysts to act with greater confidence and precision. By examining the relational connections between accounts, this method uncovers hidden subgroups of potentially collusive actors and strengthens the overall detection accuracy. It ensures that the system not only flags individual transactions but also reveals the structure of laundering networks, which is a key step for investigative and regulatory follow-up. 

We further provide the theoretical insights of using PPR for grouping money laundering behaviors. Specifically, in Theorem~\ref{thm:detectability}, we come up with a detectability condition for identifying a money-laundering group using PPR. Intuitively, PPR captures money-laundering behavior because it propagates the suspiciousness signal from a known malicious seed node through the transaction network. Nodes densely connected to the seed, likely participating in group laundering, accumulate higher PPR scores, while weakly connected nodes receive lower scores. The theorem formalizes this by showing that a well-separated group will stand out in the PPR ranking. Based on this, we modify the transition matrix in PPR by assigning higher weights to edges more likely associated with laundering transactions, which further amplifies scores for nodes strongly connected via suspicious interactions. Specifically, weighted edges increase the transition probability along suspicious links, causing PPR to concentrate more on nodes involved in laundering.

This paper demonstrates the strategic value of the proposed Cross-bank Personalized PageRank (PPR) framework. By enabling institutions to integrate Cross-bank intelligence into the ranking process, Cross-bank PPR consistently identifies substantially more true malicious accounts while flagging fewer accounts overall, compared with conventional PPR. This outcome reflects a more conservative yet more precise detection strategy, reducing unnecessary investigations while elevating the identification of genuinely high-risk entities. This paper also reveals that the incorporation of the cross-bank suspiciousness signal markedly improves precision and reduces false positives by itself, whereas the edge-level maliciousness score degrades performance, when used alone to reweight transitions. We attribute it to the reliable prediction by the detection module. However, when the edge score is fused with the cross-bank signal inside Cross-bank PPR, its weaknesses are corrected and it meaningfully enhances detection. Together, these components significantly reduce false positives, increase the concentration of malicious accounts within discovered clusters, and improve detection consistency across countries. These findings highlight that cross-institutional signal sharing, when designed with strong privacy safeguards, can meaningfully strengthen financial crime detection beyond what any single institution can achieve in isolation.

Our analysis of interpretable group identification using cross-bank Personalized PageRank reveals that money-laundering behaviors often exhibit diverse structural patterns, including fan-out, loop, gather–scatter, and hybrid mechanisms that integrate multiple strategies. The proposed cross-bank PPR method effectively uncovers the pivotal accounts orchestrating these operations, even when illicit activities are distributed across institutions and national boundaries. In particular, our approach successfully reconstructs the complete illicit group behaviors for many laundering schemes. As laundering strategies become increasingly sophisticated by mixing multiple mechanisms and fragmenting transaction flows across borders, our method continues to accurately flag primary dispersal accounts and expose partial cross-border linkages, offering critical insights into the underlying operational structure. These findings underscore that analyzing relational patterns across merged clusters not only captures the organizational essence of money-laundering schemes but also enhances interpretability, providing actionable signals for risk assessment and regulatory intervention.


% Our analysis of interpretable group identification by cross-bank Personalized PageRank reveals that group money-laundering behaviors exhibit diverse structural patterns, including fan-out, loop, gather-scatter, and hybrid mechanisms that combine multiple strategies. The cross-bank Personalized PageRank method effectively identifies the pivotal accounts controlling these operations, even when illicit activities are fragmented across institutions or countries. Specifically, our method can identify the complete group illicit behaviors for many money-laundering scheme. When the money-laundering behaviors become more and more complicated by mixing multiple laundering schemes together and fragmenting the entire process via cross-border flow, our method can still accurately flags primary dispersal accounts and reveal partial cross-border flows, which provides critical insight into the overall operation. These findings underscore that analyzing relational patterns across merged clusters not only captures the core of money-laundering schemes but also enhances interpretability, offering actionable signals for risk detection and regulatory intervention.

Finally, we propose a hierarchical reinforcement learning strategy to move from detection to action. Rather than treating AML as a static classification problem, this approach dynamically learns how to respond to suspicious activity. It considers factors such as network risk concentration, and operational costs to determine optimal intervention strategies: when to monitor, and when to escalate a freeze intervention. By balancing accuracy with cost and compliance priorities, this decision-making layer bridges the gap between automated detection and effective financial risk management. Our empirical evaluation shows that this decision-making module delivers substantial economic value. While a fixed-threshold policy prevents 50\% of illicit losses and exhibits extremely high Type II error rates, the hierarchical module prevents around 80\% of potential financial loss and dramatically reduces missed detections. Moreover, the learned thresholds vary across institutions in different countries, highlighting that optimal intervention policies are inherently context-dependent rather than universal. These findings underscore that effective AML risk management requires not just accurate prediction, but a policy that intelligently maps risk signals to timely and context-sensitive actions.

The remainder of this paper is organized as follows. We review the related works in Section 2 and Section 3 presents the proposed graph-based federated detection module, which enables privacy-preserving collaboration across institutions. Section 4 introduces the personalized PageRank and label propagation mechanisms used to uncover subgroup structures and refine detection accuracy. Section 5 details the hierarchical reinforcement learning framework that supports adaptive and coordinated decision-making. Section 6 presents the empirical results. Finally, Section 7 concludes the paper and discusses future research directions.